\documentclass[11pt]{article}
\usepackage{amsmath,amssymb,amsthm,mathtools}
\usepackage[margin=1in]{geometry}

\newtheorem{theorem}{Theorem}
\newtheorem{proposition}[theorem]{Proposition}
\newtheorem{corollary}[theorem]{Corollary}
\theoremstyle{definition}
\newtheorem{remark}[theorem]{Remark}

\newcommand{\kk}{\Bbbk}
\newcommand{\A}{\mathcal{A}}
\newcommand{\Std}{\operatorname{Std}}
\newcommand{\lm}{\operatorname{lm}}
\newcommand{\lc}{\operatorname{lc}}
\newcommand{\inn}{\operatorname{in}}
\newcommand{\pp}{\mathfrak{p}}
\DeclareMathOperator{\Spec}{Spec}

\title{Recovering a $\mathbb{Q}(q,t)$-basis from specializations:\\ what the reconstruction proves}
\author{}
\date{}

\begin{document}
\maketitle

\noindent
We reconstruct the Gr\"obner basis of a parametric ideal over $K=\mathbb{Q}(q,t)$ by solving
the ideal at many integer points $(q_0,t_0)$ and interpolating. This note states what that
procedure does and does not prove. The two theorems are standard; the point is to separate the
\emph{search} (interpolation, heuristic) from the \emph{proof} (a check with $q,t$ transcendental),
and to say precisely where the objection ``a discrete grid may sit on the bad locus'' bites.

\section{Setup}

Let $\kk=\mathbb{Q}$, $R=\kk[q,t]$, $K=\kk(q,t)$. Fix an alphabet
$X=\{x_1,\dots,x_{n-1},y_1,\dots,y_{n-1}\}$ and an admissible well-order $\prec$ on the free monoid
$\langle X\rangle$ (compatible with two-sided multiplication). Let $f_1,\dots,f_m\in R\langle X\rangle$
be the presentation (braid, inverse, zipper and untwisting relations; coefficients in
$\mathbb{Z}[q,t]$), $I={}_K(f_1,\dots,f_m)_K$ the two-sided ideal they generate in $K\langle X\rangle$,
and $\A=K\langle X\rangle/I$. We assume $\dim_K\A<\infty$, which holds in our case.

For a point $\pp=(q_0,t_0)\in\mathbb{A}^2_\kk(\kk)=\kk^2$ let $\sigma_\pp\colon R\langle X\rangle\to
\kk\langle X\rangle$ be the specialization $q\mapsto q_0,\ t\mapsto t_0$, put
$I_\pp={}_\kk(\sigma_\pp f_1,\dots,\sigma_\pp f_m)$ and $\A_\pp=\kk\langle X\rangle/I_\pp$.

For an ideal $J$ (over $K$ or $\kk$), noncommutative Buchberger--Mora reduction w.r.t.\ $\prec$
terminates when $J$ is co-finite, producing the unique reduced Gr\"obner basis $G(J)$ (monic leading
words, tails supported on standard monomials). Write $G=G(I)\subseteq K\langle X\rangle$ and
$G_\pp=G(I_\pp)\subseteq\kk\langle X\rangle$. The set of standard monomials
$\Std(J)=\{\,w\in\langle X\rangle : \lm(g)\nmid w\ \forall g\in G(J)\,\}$ is a basis of the quotient.
Each $g\in G$ has the form
\begin{equation}\label{eq:shape}
  g \;=\; m_g \;+\; \sum_{w\in T_g} c_{g,w}\,w,\qquad c_{g,w}\in K,
\end{equation}
with $m_g=\lm(g)$ and $T_g\subset\Std(I)$ finite. \emph{Reconstruction} means: recover the
coefficient functions $c_{g,w}\in K$ from their values $c_{g,w}(\pp)$ at grid points.

\section{Specialization commutes with the Gr\"obner basis off a hypersurface}

\begin{theorem}\label{thm:spec}
There is a nonzero $d\in R$ such that for every $\pp$ with $d(\pp)\neq0$:
\begin{enumerate}
\item[\textup{(i)}] every $c_{g,w}$ lies in the localization $R_d=R[d^{-1}]$, so $\sigma_\pp g$ is defined;
\item[\textup{(ii)}] $\sigma_\pp G:=\{\sigma_\pp g:g\in G\}$ is the reduced Gr\"obner basis of $I_\pp$;
\item[\textup{(iii)}] $\lm(\sigma_\pp g)=m_g$, $\ \Std(I_\pp)=\Std(I)$, and $\dim_\kk\A_\pp=\dim_K\A$.
\end{enumerate}
\end{theorem}

\begin{proof}
Run Buchberger--Mora on $\{f_i\}$ over $K$; it terminates, giving $G$, using finitely many
operations of $K$. Let $d\in R$ be a common denominator of all $c_{g,w}$ and of every leading
coefficient that was inverted (used as a pivot) during the run. Then $G\subseteq R_d\langle X\rangle$,
which is (i).

Two families of witnesses are produced by (or extractable from) the run, each a \emph{finite}
identity with coefficients in $R_d$:
\[
  \textbf{(R)}\quad f_i=\sum_{j}p_{ij}\,g_{k_j}\,q_{ij}\ \ (\text{$f_i$ reduces to }0\bmod G),
  \qquad
  \textbf{(S)}\quad S(g,g')=\sum_{j}p'_{j}\,g_{k'_j}\,q'_{j}
\]
for every $f_i$ and every overlap ambiguity $S(g,g')$ of $G$; here $p,q$ are words. Because
$d(\pp)\neq0$, $\sigma_\pp$ restricts to a ring homomorphism $R_d\to\kk$, so it carries (R) and (S)
to the same identities over $\kk$ for $\sigma_\pp G$. Thus every $\sigma_\pp f_i$ reduces to $0$
modulo $\sigma_\pp G$, and every S-polynomial of $\sigma_\pp G$ reduces to $0$ modulo $\sigma_\pp G$.
The leading words are unchanged (the $m_g$ have coefficient $1$), so Buchberger's criterion applies:
$\sigma_\pp G$ is a Gr\"obner basis of $(\sigma_\pp G)$. From (R), $I_\pp\subseteq(\sigma_\pp G)$;
from a cofactor expression $g=\sum p\,f_{i}\,q$ over $R_d$ (each $g\in I$), $\sigma_\pp g\in I_\pp$, so
$(\sigma_\pp G)\subseteq I_\pp$. Hence $(\sigma_\pp G)=I_\pp$ and $\sigma_\pp G$ is a Gr\"obner basis
of $I_\pp$ with leading words $\{m_g\}$; this gives $\Std(I_\pp)=\Std(I)$ and the equal dimensions.
Monic leading terms and standard tails are preserved by $\sigma_\pp$, so $\sigma_\pp G$ is reduced;
uniqueness of the reduced Gr\"obner basis gives $\sigma_\pp G=G_\pp$. This is (ii), (iii).
\end{proof}

Write $U=D(d)=\{\,d\neq0\,\}$ (a dense open subset of $\mathbb{A}^2$) and $B=V(d)=\mathbb{A}^2\setminus U$.

\begin{corollary}\label{cor:values}
For $\pp\in U$, the reduced Gr\"obner basis returned at $\pp$ carries, in slot $(g,w)$ of
\eqref{eq:shape}, the value $c_{g,w}(\pp)$. On $U$ the numbers we compute are point-values of the
generic coefficient functions; recovering $c_{g,w}$ from such values is interpolation.
\end{corollary}

\section{Why sampling alone cannot prove the recovery}

$B=V(d)$ is a proper Zariski-closed subset, hence a curve (a nonzero $d$ cuts out something of
codimension $\geq1$). This does \emph{not} make $B(\kk)$ finite: any component that is a rational
curve --- e.g.\ a line $\{q=1\}$ or $\{t+q^2=0\}$ --- carries infinitely many rational points.

\begin{proposition}\label{prop:gap}
No function of the samples $\{G_\pp:\pp\in\Gamma\}$ on a finite grid $\Gamma\subset\kk^2$ certifies
$\Gamma\subseteq U$. In particular, constancy of $\dim_\kk\A_\pp$ across $\Gamma$ does not.
\end{proposition}

\begin{proof}
The map $\pp\mapsto\dim_\kk\A_\pp=\#\Std(I_\pp)$ is upper semicontinuous: for a co-finite family it
satisfies $\dim_\kk\A_\pp\ge\dim_K\A$ with equality exactly on a dense open set containing $U$
(Theorem~\ref{thm:spec}(iii); semicontinuity of fibre dimension). Its level sets are locally closed,
and a level set on which the value exceeds $\dim_K\A$ can be a positive-dimensional subvariety of
$B$. A grid $\Gamma$ contained in such a level set is consistent (constant dimension, constant
leading-word set) yet lies entirely in $B$, and then the interpolant of
Corollary~\ref{cor:values} reconstructs the coefficient functions of the \emph{restricted} family,
not of $G$. Nothing intrinsic to $\{G_\pp:\pp\in\Gamma\}$ distinguishes this case.
\end{proof}

Thus interpolation is a heuristic search: it returns a candidate $\widetilde G$, correct
if $\Gamma\subseteq U$, but the samples cannot warrant $\Gamma\subseteq U$.

\section{What proves the recovery: a check with $q,t$ transcendental}

The proof is decoupled from the sampling. It runs over $K$, so the bad locus is irrelevant.

\begin{theorem}\label{thm:certify}
Let $\widetilde G=\{\tilde g_1,\dots,\tilde g_N\}\subset K\langle X\rangle$ be finite, with monic,
pairwise-nondividing leading words $m_\ell=\lm(\tilde g_\ell)$ and tails supported on
$\{w:w\prec m_\ell,\ m_{\ell'}\nmid w\ \forall\ell'\}$. Suppose, over $K$:
\begin{enumerate}
\item[\textup{(a)}] each $f_i$ reduces to $0$ modulo $\widetilde G$;
\item[\textup{(b)}] every overlap S-polynomial of $\widetilde G$ reduces to $0$ modulo $\widetilde G$;
\item[\textup{(c)}] each $\tilde g_\ell\in I$, exhibited by a two-sided cofactor
      $\tilde g_\ell=\sum_j p_j\,f_{i_j}\,q_j$.
\end{enumerate}
Then $\widetilde G=G$.
\end{theorem}

\begin{proof}
By (b) and Bergman's diamond lemma / Buchberger's criterion, $\widetilde G$ is a Gr\"obner basis of
$J:=(\widetilde G)$. By (c), $J\subseteq I$; by (a), $I\subseteq J$. Hence $J=I$ and $\widetilde G$
is a Gr\"obner basis of $I$; it is monic with standard tails, hence reduced; by uniqueness
$\widetilde G=G$.
\end{proof}

Conditions (a)--(c) contain no specialization: they are identities in $K\langle X\rangle$ with $q,t$
transcendental, so Proposition~\ref{prop:gap} does not apply to them. They are finite (normal-form
computations and cofactor expansions). Interpolation supplies the candidate $\widetilde G$;
Theorem~\ref{thm:certify} certifies it. This is the whole logic: \emph{search by specialization,
prove by transcendence.}

\begin{remark}[verifying (a)--(c) by evaluation, without symbolic algebra]\label{rem:pit}
Each reduction in (a),(b) and each cofactor in (c) is, after clearing the common denominator
$d$, a polynomial identity $P=0$ in $R\langle X\rangle$: word by word, a coefficient identity in
$R=\mathbb{Z}[q,t]$ of some bidegree $(D_q,D_t)$ \emph{read off the cleared certificate}. A
polynomial in $R$ of bidegree $(D_q,D_t)$ vanishing on a $(D_q{+}1)\times(D_t{+}1)$ integer product
grid is identically zero (iterated univariate interpolation / finite differences). So (a)--(c) are
verified by exact evaluation on such a grid --- deterministic, no probability, no genericity. This
is again a use of points, but now the points check an identity of \emph{known} degree, so the finite
grid \emph{does} suffice; contrast Proposition~\ref{prop:gap}, where the unknown was a rational
function of unbounded a priori degree tied to the unknown locus $B$.
\end{remark}

\section{Atoms}

By construction $d$ is a product of pivots, each a nonzero element of $R$. Let $a_1,\dots,a_r$ be
the distinct irreducible factors of $d$ in $R$ --- the \emph{atoms}.

\begin{proposition}\label{prop:atoms}
\leavevmode
\begin{enumerate}
\item[\textup{(i)}] Every coefficient satisfies $c_{g,w}\in R[a_1^{-1},\dots,a_r^{-1}]$, i.e.\
$c_{g,w}=N_{g,w}/\prod_i a_i^{e_i}$ with $N_{g,w}\in R$. Clearing $\prod_i a_i^{e_i}$ turns the
recovery of $c_{g,w}$ into polynomial interpolation of $N_{g,w}$.
\item[\textup{(ii)}] $B=V(d)=\bigcup_i V(a_i)$, and the locus where $\dim_\kk\A_\pp>\dim_K\A$ is
contained in $B$. The atoms index the components of the degeneration locus.
\item[\textup{(iii)}] Any $K$-valued invariant built from $\A$ by the same field operations lies in
$R[d^{-1}]$ after enlarging $d$ by finitely many further pivots; e.g.\ the Gram determinant
$\det\Theta$ of the regular trace form satisfies $\det\Theta\in R[d^{-1}]$, so for $\pp\in U$ the
algebra $\A_\pp$ is semisimple iff $\det\Theta(\pp)\neq0$, and the non-semisimple locus in $U$ is
cut out by the extra irreducible factors of $\det\Theta$ --- again atoms.
\end{enumerate}
\end{proposition}

This is the single mechanism behind the recurring factors: the atoms are the irreducible components
of the discriminantal hypersurface $B$ carrying \emph{all} denominators produced by the same
elimination --- the Gr\"obner tails, the $X_3$-cofactor denominator $t^4(t+q^2)$, and the
semisimplicity locus.

\section{A rigorous sampling route, and why we do not use it}

Proposition~\ref{prop:gap} is an obstruction only because $\deg d$ is unknown. Grant an a priori
bound $\deg d\le\Delta$ and a numerator bound $\deg N_{g,w}\le\delta$. A component $V(a_i)$ meets a
product grid $\Gamma=Q\times T$ in at most $(\deg a_i)\,|T|$ points (a plane curve of degree $e$
meets $|T|$ vertical lines in $\le e|T|$ points), so $|\Gamma\cap B|\le\Delta\,|T|$. A rational
recovery of Berlekamp--Welch type over $\Gamma$ then returns the true $N_{g,w}/\prod a_i^{e_i}$ and
locates the bad points, provided $|\Gamma|$ exceeds the interpolation requirement plus twice the
error bound. This \emph{is} a proof from samples. It fails in practice only because the general
bounds on $\Delta$ (through the sizes and coefficient degrees of the Buchberger matrices) are far too
large to make $|\Gamma|$ feasible. Hence we certify a posteriori by Theorem~\ref{thm:certify} rather
than a priori by error correction: the candidate is cheap to guess and cheap to verify, whereas the
guaranteed-correct interpolation would need an infeasible grid.

\section*{Summary}

Theorem~\ref{thm:spec} identifies what the samples are: on $U=D(d)$ they are point-values of the
generic Gr\"obner coefficients, and recovery is interpolation. Proposition~\ref{prop:gap} is the
honest limitation: a finite grid cannot certify $\Gamma\subseteq U$, so interpolation alone proves
nothing. Theorem~\ref{thm:certify} supplies the proof, over $K$ with $q,t$ transcendental, where the
bad locus has no meaning; Remark~\ref{rem:pit} makes that check a finite exact evaluation.
Proposition~\ref{prop:atoms} explains the denominators and ties them to the degeneration and
semisimplicity loci. Reconstruction is search by specialization, proof by transcendence.

\section*{Notes and references}

That a Gr\"obner basis over the rational-function field in the parameters specializes to a Gr\"obner
basis of the specialized ideal, provided a distinguished polynomial does not vanish at the
specialization point, is due to Gianni~\cite{Gianni} and, in the leading-term-stability form used in
Theorem~\ref{thm:spec}, to Kalkbrener~\cite{Kalkbrener}. The systematic parametric theory is
Weispfenning's comprehensive Gr\"obner bases~\cite{Weispfenning} and the canonical Gr\"obner cover of
Montes and Wibmer~\cite{MontesWibmer}, which partitions parameter space into locally closed strata
carrying fixed reduced-basis data; the good locus $U$ of Theorem~\ref{thm:spec} is the open stratum,
and $B$ its complement. Becker and Weispfenning~\cite{BeckerWeispfenning} is the standard text.

The algebras here are quotients of a free associative algebra, so ``Gr\"obner basis'' is meant in the
two-sided noncommutative sense. Uniqueness of the reduced basis and the confluence criterion by
overlap (``ambiguity'') S-polynomials, used in Theorems~\ref{thm:spec} and~\ref{thm:certify}, are
Bergman's diamond lemma~\cite{Bergman} and Mora's noncommutative Buchberger procedure
\cite{Mora86,Mora94}; see also Green~\cite{Green} and Ufnarovski~\cite{Ufnarovski}.

Upper semicontinuity of the fibre dimension in Proposition~\ref{prop:gap} is standard: Hartshorne
\cite{Hartshorne} (semicontinuity of fibre dimension), Eisenbud~\cite{Eisenbud} (flatness and the
Hilbert function), Matsumura~\cite{Matsumura}.

Recovery of a coefficient from its values (Corollary~\ref{cor:values}, Remark~\ref{rem:pit}) is
multivariate interpolation: dense and sparse polynomial interpolation~\cite{Zippel,BenOrTiwari},
black-box separation of numerator and denominator~\cite{KaltofenTrager}, multivariate rational
interpolation~\cite{CuytLee}, and classical Thiele/reciprocal-difference rational interpolation
\cite[\S2.2]{StoerBulirsch}. The exact-$\mathbb{Q}$ value at each grid point is itself reconstructed
from modular images --- the inner layer --- by modular and traced Gr\"obner computation
\cite{Traverso,Arnold,NoroYokoyama} with rational-number reconstruction~\cite{WangGuyDavenport}.

The error-correcting recovery of \S7 is the decoding view of interpolation: the Welch--Berlekamp
algorithm~\cite{WelchBerlekamp,GemmellSudan}, its list-decoding extension~\cite{GuruswamiSudan}, and
sparse-interpolation codes decoding beyond half the minimum distance~\cite{KaltofenPernet}.

\begin{thebibliography}{99}

\bibitem{Weispfenning} V.~Weispfenning, \emph{Comprehensive Gr\"obner bases},
  J. Symbolic Comput. \textbf{14} (1992), no.~1, 1--30. DOI 10.1016/0747-7171(92)90023-W.
\bibitem{Kalkbrener} M.~Kalkbrener, \emph{On the stability of Gr\"obner bases under specializations},
  J. Symbolic Comput. \textbf{24} (1997), no.~1, 51--58. DOI 10.1006/jsco.1997.0113.
\bibitem{Gianni} P.~Gianni, \emph{Properties of Gr\"obner bases under specializations},
  EUROCAL~'87 (J.~H.~Davenport, ed.), Lecture Notes in Comput. Sci. \textbf{378},
  Springer, 1989, 293--297. DOI 10.1007/3-540-51517-8\_128.
\bibitem{MontesWibmer} A.~Montes and M.~Wibmer, \emph{Gr\"obner bases for polynomial systems with
  parameters}, J. Symbolic Comput. \textbf{45} (2010), no.~12, 1391--1425.
  DOI 10.1016/j.jsc.2010.06.017.
\bibitem{BeckerWeispfenning} T.~Becker and V.~Weispfenning (with H.~Kredel), \emph{Gr\"obner Bases:
  A Computational Approach to Commutative Algebra}, Grad. Texts in Math. \textbf{141},
  Springer, 1993.

\bibitem{Bergman} G.~M.~Bergman, \emph{The diamond lemma for ring theory},
  Adv. Math. \textbf{29} (1978), no.~2, 178--218. DOI 10.1016/0001-8708(78)90010-5.
\bibitem{Mora86} T.~Mora, \emph{Gr\"obner bases for non-commutative polynomial rings},
  AAECC-3 (Grenoble, 1985), Lecture Notes in Comput. Sci. \textbf{229}, Springer, 1986, 353--362.
  DOI 10.1007/3-540-16776-5\_740.
\bibitem{Mora94} T.~Mora, \emph{An introduction to commutative and noncommutative Gr\"obner bases},
  Theoret. Comput. Sci. \textbf{134} (1994), no.~1, 131--173. DOI 10.1016/0304-3975(94)90283-6.
\bibitem{Green} E.~L.~Green, \emph{Noncommutative Gr\"obner bases, and projective resolutions},
  Computational Methods for Representations of Groups and Algebras (Essen, 1997),
  Progr. Math. \textbf{173}, Birkh\"auser, 1999, 29--60. DOI 10.1007/978-3-0348-8716-8\_2.
\bibitem{Ufnarovski} V.~Ufnarovski, \emph{Introduction to noncommutative Gr\"obner bases theory},
  Gr\"obner Bases and Applications (Linz, 1998), London Math. Soc. Lecture Note Ser. \textbf{251},
  Cambridge Univ. Press, 1998, 259--280. DOI 10.1017/CBO9780511565847.015.

\bibitem{Hartshorne} R.~Hartshorne, \emph{Algebraic Geometry}, Grad. Texts in Math. \textbf{52},
  Springer, 1977.
\bibitem{Eisenbud} D.~Eisenbud, \emph{Commutative Algebra with a View Toward Algebraic Geometry},
  Grad. Texts in Math. \textbf{150}, Springer, 1995.
\bibitem{Matsumura} H.~Matsumura, \emph{Commutative Ring Theory} (trans. M.~Reid),
  Cambridge Stud. Adv. Math. \textbf{8}, Cambridge Univ. Press, 1989.

\bibitem{Zippel} R.~Zippel, \emph{Interpolating polynomials from their values},
  J. Symbolic Comput. \textbf{9} (1990), no.~3, 375--403. DOI 10.1016/S0747-7171(08)80018-1.
\bibitem{BenOrTiwari} M.~Ben-Or and P.~Tiwari, \emph{A deterministic algorithm for sparse
  multivariate polynomial interpolation}, STOC~'88, ACM, 1988, 301--309. DOI 10.1145/62212.62241.
\bibitem{KaltofenTrager} E.~Kaltofen and B.~M.~Trager, \emph{Computing with polynomials given by
  black boxes for their evaluations: greatest common divisors, factorization, separation of
  numerators and denominators}, J. Symbolic Comput. \textbf{9} (1990), no.~3, 301--320.
  DOI 10.1016/S0747-7171(08)80015-6.
\bibitem{CuytLee} A.~Cuyt and W.-s.~Lee, \emph{Sparse interpolation of multivariate rational
  functions}, Theoret. Comput. Sci. \textbf{412} (2011), no.~16, 1445--1456.
  DOI 10.1016/j.tcs.2010.11.050.
\bibitem{StoerBulirsch} J.~Stoer and R.~Bulirsch, \emph{Introduction to Numerical Analysis},
  3rd ed., Texts in Appl. Math. \textbf{12}, Springer, 2002.

\bibitem{WelchBerlekamp} L.~R.~Welch and E.~R.~Berlekamp, \emph{Error correction for algebraic block
  codes}, U.S. Patent 4,633,470, 1986.
\bibitem{GemmellSudan} P.~Gemmell and M.~Sudan, \emph{Highly resilient correctors for polynomials},
  Inform. Process. Lett. \textbf{43} (1992), no.~4, 169--174. DOI 10.1016/0020-0190(92)90195-2.
\bibitem{GuruswamiSudan} V.~Guruswami and M.~Sudan, \emph{Improved decoding of Reed--Solomon and
  algebraic-geometry codes}, IEEE Trans. Inform. Theory \textbf{45} (1999), no.~6, 1757--1767.
  DOI 10.1109/18.782097.
\bibitem{KaltofenPernet} E.~L.~Kaltofen and C.~Pernet, \emph{Sparse polynomial interpolation codes
  and their decoding beyond half the minimum distance}, ISSAC~2014, ACM, 2014, 272--279.
  DOI 10.1145/2608628.2608660.

\bibitem{Traverso} C.~Traverso, \emph{Gr\"obner trace algorithms}, ISSAC~'88 (Rome, 1988),
  Lecture Notes in Comput. Sci. \textbf{358}, Springer, 1989, 125--138.
  DOI 10.1007/3-540-51084-2\_12.
\bibitem{Arnold} E.~A.~Arnold, \emph{Modular algorithms for computing Gr\"obner bases},
  J. Symbolic Comput. \textbf{35} (2003), no.~4, 403--419. DOI 10.1016/S0747-7171(02)00140-2.
\bibitem{NoroYokoyama} M.~Noro and K.~Yokoyama, \emph{Usage of modular techniques for efficient
  computation of ideal operations}, Math. Comput. Sci. \textbf{12} (2018), no.~1, 1--32.
  DOI 10.1007/s11786-017-0325-1.
\bibitem{WangGuyDavenport} P.~S.~Wang, M.~J.~T.~Guy, and J.~H.~Davenport, \emph{$p$-adic
  reconstruction of rational numbers}, ACM SIGSAM Bull. \textbf{16} (1982), no.~2, 2--3.
  DOI 10.1145/1089292.1089293.

\end{thebibliography}

\end{document}
